% INTRODUCTION. Numbers ONLY via \R... macros from generated/numbers.tex.
An LLM agent's prompt is assembled, not written. Before each decision a harness collects tool schemas, retrieved
documents, worked examples, notes and conversation history, and must fit them into a token budget set by latency, cost or
the model's window~\cite{cesurvey2025}. What to keep is usually decided by heuristics: a relevance score against the
goal, a coverage objective over candidate passages~\cite{pacms2026}, a compressor that deletes tokens~\cite{llmlingua2_2024},
or fixed priorities written by an engineer. None of these asks the question a budgeted system actually faces:
\emph{how much does this kind of context change task success for this model on this task family?} The answer is not
obvious. Models use long contexts unevenly~\cite{lostmiddle2024}, and a block can be relevant to the goal yet add
nothing the model needs, or be redundant with another block that is already included.

This paper asks: when context consists of heterogeneous types, can we measure how much each type contributes to task
success, and can those measurements be reused for budgeted assembly? We separate three levels of the problem
(Section~\ref{sec:problem}): \emph{type value} (which kinds of context matter for a model and task family),
\emph{block relevance} (which particular blocks matter for an instance), and \emph{budget allocation} (what to send under a
token limit, and when to stop). A \emph{marginal context value (MCV) profile} answers the first level. For a (model, task
family) pair it records the measured drop in success when every block of a type is removed, the value of the type's
short variant, and pairwise interactions between types, each with a bootstrap interval. Profiles are estimated once,
offline, by counterfactual ablation on development instances, with a length-neutral control. For the second level we
train a block-value model on leave-one-block-out deltas from the same runs. For the third, a CP-SAT compiler maximizes
predicted value under the budget, dependencies and variants without calling the language model, and reports a
certificate of optimality for that predicted objective.

Prior work values context in related ways but for other decisions: attribution methods explain one generation by
ablating its sources~\cite{contextcite2024,attribot2024}, Shapley-style methods value demonstrations, prompts, data
points or skills~\cite{demoshapley2024,promptshapley2023,datashapley2019,skillsv2026}, and recent context-allocation
studies measure the causal effect of context in generative search~\cite{ctxalloc2026}. Section~\ref{sec:related}
positions MCV profiles against these and against budgeted selectors. The distinguishing combination is a reusable,
type-level, interval-reported value table measured per model and task family and compiled under a hard token budget.

We evaluate the approach under a pre-registered protocol (hypotheses, margins, sample sizes and the method version fixed
before any test run, recorded as a commit in the released repository; it was not deposited with an external registry) on three task families: multi-hop question answering over HotpotQA documents with distractors,
tool selection and argument filling over a \RNumTools-tool catalog built from BFCL, and recall of the latest value of a
user attribute from \RNumSessions synthetic sessions. All models run locally (Qwen3-4B-Instruct for the main study,
Qwen3-8B for transfer). We compare against six baselines, including LLMLingua-2 and a PACMS-style coverage selector.

The primary result is negative, and we lead with it. \textbf{MCV did not beat relevance at block selection:}
MCV-compiled context did not significantly exceed embedding relevance (the best baseline) in area under the
success--budget curve in any family, and was \RHOneHistoryAbs{} lower on history. The cause is at the block level. Single
deletions rarely change the outcome, so the block-value model learned proxies such as recency and length. At the smallest
history budget it kept the session holding the answer in \RRecallOurs of plans, against \RRecallRel for relevance.
\textbf{Efficiency held, but is not specific to MCV.} The efficiency hypothesis, added after development pilots but before
any test run, held in \RHOnebPass of three families: at an 8k budget, which full contexts never reach, MCV pruned
\RTokCutTool of prompt tokens on tool use and \RTokCutHistory on history with non-inferior success. A pre-specified
post-hoc follow-up showed that relevance with a calibrated stopping rule pruned more (\RTokCutStopTool on tool use, also
non-inferior), and that a hybrid in which the profile admits types and relevance ranks blocks did not beat relevance.
Interaction terms had no detectable effect. \textbf{As a diagnostic, the profiles were informative.} Tool schemas,
documents and history carried most of the value (\RMcvToolToolSchema, \RMcvFactsFact, \RMcvHistoryHistoryTurn),
worked examples little on their own, and shortened tool schemas lost most of theirs. History and the profile note
substituted for each other (interaction \RIntHistoryHistoryTurnNote), and the length-neutral control moved no value by
more than \RPadDiffMax. Qwen3-4B and Qwen3-8B type values were rank-correlated ($\rho=\RRho$), but outcome agreement under
transferred and native profiles may reflect compiler insensitivity as much as transfer.

\textbf{Contributions.}
\begin{itemize}
\item \emph{A three-level view of context assembly} (type value, block relevance, budget allocation), with evidence that
the levels should not be served by one signal: measured type value was informative, while block selection was best served
by embedding relevance (Sections~\ref{sec:problem} and~\ref{sec:eval}).
\item \emph{MCV profiles as a reusable type-level diagnostic:} an estimator of per-(model, family) type values,
short-variant values and interactions with shrinkage, bootstrap intervals and a length-neutral control, with measured
stability (ten to twenty development instances in subsampling), cost and Qwen3-4B-to-8B agreement.
\item \emph{A controlled, pre-registered evaluation} of ten context policies on three task families, including a failed
primary hypothesis, a pre-specified follow-up, and evidence that leave-one-block-out deltas were too sparse to train a
block ranker that improves on embedding similarity.
\item \emph{A dependency- and variant-aware exact budget compiler} whose certificate is bounded to the predicted
objective, not task success (Section~\ref{sec:design}).
\item \emph{An open artifact:} code, data adapters, logs of every model call, and scripts that regenerate every number,
table and figure.
\end{itemize}
